\chapter{Návrh informačního systému}
\label{ch:design}
Tato kapitola navazuje na předchozí analýzy a~legislativní rozbory a~definuje konkrétní technickou architekturu informačního systému pro turistický kemp Olšovec.
Cílem je popsat řešení, jež zajistí spolehlivost, bezpečnost a~dostupnost systému při zachování provozní pružnosti nezbytné pro sezónní charakter areálu.
Návrh vymezuje rozdělení systému na komponenty a~jejich odpovědnosti, principy komunikace mezi nimi, datový model vycházející z~dříve popsaného doménového modelu a~zásadní části uživatelského rozhraní.
Důraz je kladen rovněž na integraci s~externími službami a~na propojení s~fyzickou infrastrukturou areálu, jejíž součástí jsou vjezdové závory ovládané uzavřeným proprietárním systémem.

\section{Komponenty systému a~jejich role}
\label{sec:komponenty}
Tato sekce popisuje rozklad systému na jednotlivé komponenty, vymezuje jejich odpovědnosti a~zachycuje způsob jejich vzájemné spolupráce v~rámci výsledného celku.
Architektura je navržena s~důrazem na vysokou dostupnost, nízkou latenci a~minimalizaci závislosti na externí konektivitě, což odpovídá specifickým provozním podmínkám kempu.

Centrální výpočetní jednotkou systému je lokální aplikační server provozovaný v~rámci vnitřní sítě areálu.
Toto umístění je motivováno snahou zajistit kontinuitu provozu i~v~případě výpadku internetového připojení, kdy musejí být klíčové procesy\,--\,evidence ubytování, fakturace nebo tisk dokladů\,--\,plně funkční.
První z~komponent provozovaných na lokálním serveru je backendová aplikační vrstva, do níž je důsledně centralizována veškerá aplikační logika.
Vrstva představuje jediný autoritativní zdroj pravdy pro všechny klientské aplikace a~poskytuje jim přístup k~funkcionalitě prostřednictvím jednotného aplikačního rozhraní, jehož detailní návrh je rozebrán v sekci \ref{sec:rozhrani}.
Druhou komponentu na lokálním serveru tvoří perzistentní datová vrstva, která slouží výhradně jako úložiště dat.
Striktní oddělení odpovědností zajišťuje jednotné chování systému napříč všemi klientskými aplikacemi a~eliminuje riziko nekonzistence způsobené duplikací logiky.
Přístup k~datům je možný výhradně skrze backendovou vrstvu, která vystupuje jako jediný bod pro jejich čtení i~zápis.

Prezentační vrstva je tvořena trojicí interních klientských aplikací, provozovaných výhradně v~rámci vnitřní sítě areálu, a~doplněna o~veřejně dostupný rezervační portál pro externí hosty.
Hlavní administrativní rozhraní je koncipováno jako webová aplikace provozovaná na počítači recepce a~slouží jako primární pracovní nástroj personálu pro správu rezervací, hostů, fakturace i~dalších provozních procesů.
Tabletová aplikace pro hosty je provozována na samostatném zařízení umístěném u~recepčního pultu a~slouží jako rozhraní pro přímou komunikaci s~hostem během odbavovacího procesu.
Hlavní úlohou je zobrazování shrnutí účtenky a~sběr digitálních podpisů na přihlašovací tiskopisy cizince.
Oddělení tohoto rozhraní od pracovní stanice recepce přispívá k~plynulosti odbavení a~k~jasnému rozdělení informací určených pro personál a~pro hosta.
Mobilní aplikace pro personál slouží jako pracovní nástroj různých rolí pohybujících se po areálu\,--\,recepční ji využívají při kontrolách hostů a~jejich vozidel, provozní personál pro evidenci své činnosti přímo v~místě výkonu práce.
Jádrem aplikace jsou dvě funkce využívající systémová rozhraní mobilního zařízení k~rychlé identifikaci objektů v~terénu.
První z~nich je automatizované čtení registračních značek vozidel pomocí optického rozpoznávání znaků (OCR), které umožňuje kontrolu oprávněnosti vjetých vozidel bez nutnosti návratu na recepci.
Druhou je čtení QR kódů ze štítků umístěných na jednotlivých stanech, díky čemuž lze v~terénu okamžitě identifikovat konkrétní ubytovací jednotku, ověřit její obsazenost a~přiřazenou rezervaci či zaznamenat související provozní událost\,--\,například nahlášenou závadu nebo dokončený úklid.

Architekturu doplňují dvě lokálně provozované služby, jejichž rolí je zprostředkovávat komunikaci mezi softwarovou částí systému a~fyzickými zařízeními v~areálu.
Obě služby vystavují vlastní lokální aplikační rozhraní, liší se však tím, která komponenta systému je jejich volajícím\,--\,toto rozdělení vychází z~povahy operace, kterou každá z~nich obsluhuje.

První z~nich je tisková služba, jež řeší omezení webových prohlížečů v~oblasti přímého přístupu k~tiskovým zařízením.
Tuto službu volá webová aplikace recepce.
Princip jediného zdroje pravdy zůstává zachován\,--\,služba nezpracovává doménová data ani neprovádí žádné aplikační rozhodnutí; přijímá výhradně připravená tisková data z~backendu předaná přes prohlížeč a~zajišťuje jejich okamžitý tisk na připojené tiskárně.
Namísto využití tiskového dialogu prohlížeče jsou tak data doručena přímo službě, čímž dochází k~výraznému zkrácení doby odbavení hosta a~ke zvýšení ergonomie práce recepčního personálu.
Druhou službou je integrační vrstva pro vjezdovou infrastrukturu, která tvoří komunikační uzel mezi systémem a~uzavřeným proprietárním softwarem stávajícího řešení vjezdových závor.
Služba je provozována na počítači přímo připojeném k~ovládací jednotce závor a~pokyny přijímá výhradně z~backendu, čímž je zachována centralizace logiky\,--\,backend ověřuje oprávnění a~zachovává konzistenci dat.
Mechanismus integrace s~původním systémem je detailně popsán v~sekci \ref{sec:zavory}.

Celková architektura systému, viz obrázek \ref{fig:komunikace}, je navržena jako úzce integrovaný celek, v~němž jednotlivé komponenty plní jasně vymezené role a~zároveň sdílejí společné datové jádro.
Tato úroveň integrace softwarových i~hardwarových prvků vytváří základ pro efektivní, spolehlivou a~moderní správu turistického kempu.

\begin{figure}[h]
    \centering
    \includegraphics[width=1.0\textwidth]{diagram.png}
    \caption{Komponenty systému a~směry datové komunikace mezi nimi.}
    \label{fig:komunikace}
\end{figure}

\section{Publikace systému do internetu}
\label{sec:publikace}
Architektura systému je primárně koncipována pro lokální nasazení v~rámci vnitřní sítě areálu, což se z~bezpečnostního i~provozního hlediska osvědčuje pro interní agendu personálu.
Pro vytváření rezervací hosty však musí být systém selektivně dostupný také z~veřejné sítě\,--\,prostřednictvím veřejného rezervačního portálu.
Tato dvojí dostupnost si vyžaduje samostatnou publikační vrstvu, jejíž rolí je zpřístupnit externím klientům výhradně předem vymezenou podmnožinu aplikačního rozhraní určenou pro veřejnou rezervační agendu.
Administrativní a~ostatní interní rozhraní zůstávají dostupná pouze z~vnitřní sítě a~touto vrstvou neprocházejí, čímž je zachována striktní separace veřejné a~interní plochy systému a~minimalizována potenciální plocha pro útoky.

Aby návrh nebyl pevně vázán na konkrétní podmínky internetového připojení kempu, je publikační vrstva pojata abstraktně jako vstupní brána, jejíž konkrétní provedení může nabývat dvou alternativních forem.
Z~pohledu backendu a~ostatních komponent jsou obě varianty zaměnitelné\,--\,obě prezentují veřejnosti tutéž podmnožinu aplikačního rozhraní a~na úrovni interní logiky se chovají identicky.

První variantou je reverzní proxy na lokálním serveru přístupná přes pevnou veřejnou IP adresu.
Reverzní proxy v~tomto uspořádání naslouchá na portu 443 a~směruje validní požadavky na publikované endpointy backendu.
Příchozí provoz na pevnou IP adresu je k~serveru doručen přesměrováním na směrovači poskytovatele internetového připojení.
Toto řešení je vhodné v~prostředí, kde poskytovatel garantuje přidělení statické IP adresy a~umožňuje konfiguraci přesměrování příchozích spojení do lokální sítě.

Druhou variantou je reverzní tunel s~veřejně dostupným vstupním bodem.
V~tomto uspořádání místní agent navazuje trvalé odchozí spojení vůči externí tunelovací službě, která vystavuje veřejný hostname a~veškerý příchozí provoz směřovaný na tento hostname protuneluje zpět agentovi a~následně do publikační vrstvy.
Toto řešení nevyžaduje pevnou IP adresu ani otevírání příchozích portů na straně lokální sítě\,--\,příchozí provoz je z~pohledu lokálního směrovače fakticky odpovědí na odchozí spojení, terminaci TLS zajišťuje sama tunelovací služba.
Varianta je vhodná v~prostředí s~dynamickou IP adresou, sdíleným připojením za NAT na straně poskytovatele, nebo v~situaci, kdy kemp nedisponuje technickým zázemím pro správu vlastního TLS certifikátu a~konfigurace směrovače.

Tato architektonická pružnost umožňuje zvolit konkrétní variantu podle reálných síťových podmínek dané lokality a~v~budoucnu mezi nimi případně přejít bez zásahu do zbytku systému.

\section{Návrh aplikačního rozhraní}
\label{sec:rozhrani}
Tato sekce vymezuje principy, podle kterých probíhá komunikace mezi backendem a~klientskými aplikacemi systému.
Cílem není podat úplný výčet operací, nýbrž popsat pravidla, podle nichž jsou navrhovány, autentizovány a~zpřístupňovány různým kategoriím klientů.

Komunikace mezi klienty a~backendem je realizována ve stylu \texttt{REST} nad protokolem \texttt{HTTP}, s~přenosem zpráv ve formátu \texttt{JSON}.
Bezstavovost zjednodušuje obnovu po výpadcích, jednotná sémantika metod protokolu \texttt{HTTP} snižuje kognitivní zátěž při vývoji různorodých klientů a~formát \texttt{JSON} je pro všechny zamýšlené platformy dostupný bez specifických závislostí.
Alternativy jako \texttt{RPC} nebo grafová dotazovací rozhraní by úměrný přínos nenabízely\,--\,\texttt{REST} dobře odpovídá doménově orientované struktuře operací nad jasně vymezenými entitami a~při omezeném počtu klientů s~podobnými potřebami nepřináší \texttt{GraphQL} dostatečnou hodnotu, jež by ospravedlnila vyšší implementační složitost.

Zdroje rozhraní odpovídají hlavním entitám datového modelu\,--\,rezervacím, hostům, ubytovacím jednotkám, fakturám.
Adresní struktura odráží vztahy mezi entitami: cesta podřazeného zdroje navazuje na cestu nadřazené entity (například doklady svázané s~konkrétní rezervací), hlavní entity jsou zároveň dosažitelné přímo přes vlastní cestu.
Standardní operace vycházejí z~ustálené sémantiky metod protokolu \texttt{HTTP} (\texttt{GET}, \texttt{POST}, \texttt{PUT}, \texttt{PATCH}, \texttt{DELETE}), pro přehledové scénáře jsou definovány konvence dotazovacích parametrů pro filtrování a~textové vyhledávání.
Operace vyžadující atomický přechod stavu napříč více objekty\,--\,například dokončení odbavovacího procesu hosta nebo ohlášení cizinců cizinecké policii\,--\,jsou exponovány jako dedikované akční endpointy pod příslušným zdrojem ve tvaru \texttt{POST /resource/\{id\}/akce}.

Identita klienta je ověřována tokenovým mechanismem zasílaným v~hlavičkách každého požadavku.
Interní uživatelé\,--\,personál kempu\,--\,získávají po úspěšném přihlášení krátkodobý přístupový token doplněný o~delší obnovovací token, jenž umožňuje prodloužit přístupový token do určité časové hranice, po které se uživatel musí opět autentizovat.
Součástí přístupového tokenu jsou nároky popisující identitu uživatele a~jemu přidělené role; autorizace endpointů vychází z~rolí odpovídajících pracovním pozicím\,--\,recepční, účetní, provozní personál a~správce.
Veřejný rezervační portál naproti tomu nepracuje s~trvalou identitou hosta; bezpečnostní obálkou operací nad konkrétní rezervací je kryptograficky bezpečný náhodný token vygenerovaný při jejím založení a~předaný hostovi odkazem zaslaným na e-mailovou adresu.
Token nepodléhá rotaci, neboť přes portál není možné získat citlivá data nad rámec údajů, které host sám zadal při tvorbě rezervace.
Operace prohlížení dostupnosti a~založení nové rezervace jsou ze své povahy plně anonymní.

Příslušnost endpointu k~veřejné nebo interní podmnožině rozhraní je vyjádřena deklarativně v~metadatech jednotlivých operací, nikoliv prostřednictvím odlišné \texttt{URL} předpony.
Tento zápis slouží jako autoritativní zdroj seznamu veřejně přístupných cest, který přebírá zprostředkovatel připojení do internetu popsaný v~sekci \ref{sec:publikace}; do veřejné sítě tak procházejí výhradně požadavky směřující na endpointy z~tohoto seznamu, ostatní rozhraní zůstává dostupné pouze z~vnitřní sítě.
Veřejnou podmnožinu tvoří endpointy pro zjištění dostupnosti ubytování, založení a~správu vlastní rezervace, online check-in hostů a~stažení dokladu; administrativní agenda zůstává vyhrazena pro interní část.

Pro signalizaci výsledku požadavku jsou používány standardní stavové kódy protokolu \texttt{HTTP} doplněné o~strukturovanou chybovou odpověď ve formátu inspirovaném doporučením pro problémové detaily (\texttt{RFC~7807}).
Tělo chybové odpovědi obsahuje strojově čitelný kód, srozumitelný popis a~v~případě validačních problémů detailní rozpis nesouladu po jednotlivých polích, čímž jsou klientům jednoznačně rozlišeny autorizační chyby, validační chyby, konfliktní stavy a~chyby technické povahy.
Validace vstupních dat probíhá výhradně na straně backendu jakožto autoritativního zdroje pravdy.

\section{Datový model}
Datový model, konkrétně schéma relační databáze zobrazené na obrázku \ref{fig:db-diagram}, přímo vychází z~dříve popsaného doménového modelu viz obrázek \ref{fig:domain-model}, jenž obsahuje detailní popis jednotlivých entit.
Tato sekce se proto zaměřuje pouze na prvky, jež jsou pro relační realizaci specifické a~v~doménovém modelu nemají přímou paralelu.

Prvním z~nich jsou mapovací tabulky vyjadřující vztahy typu mnoho-k-mnoha, které relační databáze nedokáže reprezentovat přímo.
V~navrženém schématu plní tuto roli například \texttt{GroupReservationSpots}, \texttt{EventSpotGroupItems} či \texttt{SpotOofItems}.

Druhou skupinou prvků, které v~doménovém modelu zcela absentují, jsou tabulky identitního subsystému poskytující datovou strukturu pro uživatele systému, jejich přidělení k~pracovním rolím popsaným v~sekci \ref{sec:rozhrani}, vystavené tokeny a~údaje pro alternativní bezheslové přihlášení.
Identitní subsystém je čistě technickou vrstvou zajišťující autentizaci a~autorizaci přístupu, doménovému modelu kempu proto nenáleží a~objevuje se až v~tomto diagramu.

\begin{figure}[htb!]
    \centering
    \includegraphics[width=1.0\textwidth]{db-schema.pdf}
    \caption{Návrhový diagram schématu relační databáze\protect\footnotemark.}
    \label{fig:db-diagram}
\end{figure}
\footnotetext{Interaktivní verze viz \url{https://dbdiagram.io/d/DP-6a06fe2f9f1f8ec47b20a0cf}.}

\section{Integrace externích systémů}
Integrace s~externími rozhraními je v~navrženém systému klíčovým prostředkem k~zajištění datové integrity.
Namísto manuálního přepisování údajů systém využívá standardizovaná rozhraní k~validaci a~automatickému doplňování informací, což snižuje chybovost a~urychluje proces odbavení hostů v~exponovaných časech.

\subsection{Automatizovaná identifikace právnických osob}
Pro zjednodušení procesů spojených s~ubytováním a~fakturací právnických osob, jako jsou školy či firmy, systém integruje rozhraní Administrativního registru ekonomických subjektů (ARES)\footnote{Dokumentace viz \url{https://ares.gov.cz/stranky/vyvojar-info}.}.
Komunikace s~registrem probíhá skrze REST API, jež umožňuje přímý přístup k~veřejným datům o~ekonomických subjektech v~České republice.
Systém po zadání identifikačního čísla (IČO) odešle dotaz na toto rozhraní, jež vrací strukturovanou odpověď ve formátu \texttt{JSON} nebo \texttt{XML}.
Získané údaje jsou automaticky zpracovány a~použity pro předvyplnění fakturačních polí, což garantuje správnost dokladů v~souladu s~aktuálními údaji v~registru bez nutnosti jejich ručního opisování.

\subsection{Digitální verifikace identity fyzických osob}
V~rámci technického návrhu integrace se systémem eDoklady\footnote{Dokumentace viz \url{https://edoklady.gov.cz/docs/technicka-specifikace.zip}.} byly posouzeny čtyři nabízené způsoby realizace, které státní správa pro ověřovatele nabízí.
Prvním je využití nativní mobilní aplikace, jež je však pro potřeby komplexního informačního systému neintegrovatelná, neboť výsledná data zůstávají izolována v~mobilním zařízení a~vyžadují manuální přepis do databáze.
Další dvě cesty představuje takzvaná Webová čtečka ověřovatele, která funguje buď na principu zkopírování údajů ve formátu \texttt{JSON} do systémové schránky (vhodné pro manuální vkládání), nebo na principu automatizovaného vyzvednutí dat backendem, což je mechanismus podobný autorizačnímu protokolu OAuth.
Čtvrtou a~technologicky nejvyspělejší variantou je využití přímého Integračního API, jež zajišťuje čistou komunikaci typu server-to-server (B2B).

Pro účely informačního systému kempu bylo jako optimální zvoleno právě Integrační API.
Tato volba je nejvhodnější z~hlediska uživatelské přívětivosti i~technické implementace, jelikož umožňuje plnou automatizaci datového toku přímo do backendu bez nutnosti využívat mezikroky v~prohlížeči.
Z~hlediska vývoje se jedná o~nejčistší řešení, jež nabízí vysokou míru kontroly nad celým procesem a~je přímo určené pro systémy s~vyšší frekvencí ověřování.

Proces předání dat je v~tomto režimu navržen jako přímočarý tříkrokový cyklus, jenž minimalizuje zátěž pro hosta i~personál.
V~prvním kroku backend systému vygeneruje unikátní požadavek na ověření a~zobrazí jej ve formě QR kódu na tabletu u~recepčního pultu.
Host tento kód naskenuje svou aplikací eDoklady, čímž dojde k~načtení požadavku na přenos osobních dat.
Ve druhém kroku host ve své mobilní aplikaci eDoklady vidí přehled požadovaných údajů a~odsouhlasí jejich předání kempu.
Tímto úkonem udělí digitální souhlas k~jednorázovému poskytnutí identity.

V~závěrečné fázi backend pomocí zabezpečeného tokenu asynchronně údaje načte.
Data dorazí ve strukturovaném formátu \texttt{JSON} a~jsou následně zpracována.
Tento postup zaručuje, že do systému vstupují pouze ověřená a~digitálně podepsaná data, čímž se eliminuje riziko chyb vzniklých manuálním opisováním z~fyzických dokladů.
Kompletní proces ověření s~využitím této varianty je zobrazen na obrázku \ref{fig:edoklady-server}.

\begin{figure}[hbt!]
    \centering
    \includegraphics[width=1.0\textwidth]{edoklady.pdf}
    \caption{Sekvenční diagram načtení identity pomocí Integračního API systému eDoklady.}
    \label{fig:edoklady-server}
\end{figure}

\subsection{Inteligentní našeptávání a~validace adresních údajů}
Sběr adresních údajů není pouze administrativní volbou, ale zákonnou povinností ubytovatele, který musí vést evidenci všech ubytovaných osob pro účely hlášení místních poplatků obci a~hlášení pobytu cizinců.
Pro zajištění maximální kvality těchto dat a~eliminaci manuálních chyb je systém navržen tak, aby uživateli při psaní aktivně nabízel validní adresní body prostřednictvím externích registrů.
Tento problém částečně řeší možnost využívání digitální identity, většina hostů však stále preferuje fyzický občanský průkaz.

Pro dosažení plné správnosti adres v~rámci České republiky a~zároveň zajištění podpory pro zahraniční turisty systém využívá kombinaci dvou doplňujících se zdrojů komunikujících přes REST API.
Volba zdroje je řízena explicitním přepínačem ve formuláři (tuzemská versus zahraniční adresa); přepínač je výchozí podle zadané státní příslušnosti hosta, obsluha jej však může kdykoli změnit.
Pro adresy na území České republiky je primárním zdrojem Registr územní identifikace, adres a~nemovitostí (RÚIAN), jenž poskytuje geokódovací služby skrze REST API.
Tato integrace nabízí výhodu v~podobě neomezeného počtu dotazů, což je pro kemp s~vysokou sezónní obrátkou hostů klíčové.
Systém využívá konkrétně službu \texttt{findAddressCandidates}\footnote{Dokumentace viz \url{https://geoportal.cuzk.gov.cz/Dokumenty/Rozhrani_GeocodeSOE.pdf}.}, jež na základě textového řetězce vrací seznam nejpravděpodobnějších adresních bodů v~rámci České republiky.

Pro hosty přijíždějící ze zahraničí, u~kterých je evidence adresy rovněž povinná, systém přepíná na REST API služby Mapy.com\footnote{Dokumentace viz \url{https://developer.mapy.com/cs/rest-api/funkce/geokodovani/}.}.
Tato platforma poskytuje pokročilé našeptávání a~vyhledávání adres v~globálním měřítku.
V~rámci navrženého bezplatného tarifu\footnote{Informace o cenovém tarifu viz \url{https://developer.mapy.com/cs/cena/cenova-kalkulacka/}.} umožňuje tato integrace až 62\,500 požadavků za měsíc, což s~velkou rezervou pokrývá potřeby kempu.

\subsection{Hlášení pobytu cizince}
Pro splnění zákonné povinnosti hlášení pobytu cizinců ubytovatelem definuje Služba cizinecké policie ČR dvě technické cesty pro přenos dat.
První variantou je generování datového souboru ve formátu UNL, což je specificky strukturovaný textový dokument obsahující údaje o~ubytovaných osobách, jejž ubytovatel následně manuálně nahrává prostřednictvím webového rozhraní aplikace Ubyport.
Druhou metodou je využití automatizované webové služby\footnote{Dokumentace viz: \url{https://policie.gov.cz/clanek/informace-pro-vyvojare.aspx}} založené na protokolu SOAP (Simple Object Access Protocol), jež umožňuje přímé propojení informačního systému ubytovatele se servery Policie ČR pro automatické zasílání dat.

V~rámci návrhu tohoto informačního systému je jako výhradní metoda zvoleno Integrační API.
Důvodem pro upuštění od podpory UNL souborů je snaha o~maximální eliminaci manuálních procesů a~lidské chybovosti; manuální generování a~následné nahrávání souborů představuje administrativní zátěž, jež neodpovídá architektuře navrženého systému.
Webová služba navíc poskytuje okamžitou zpětnou vazbu ve formě takzvané Doručenky, díky které systém ihned ví, zda bylo hlášení přijato v~pořádku, nebo zda obsahuje chyby, jež je nutné opravit.
V~případě, že Doručenka oznámí neúspěch, je hlášení zařazeno do fronty pro opravu a~personál recepce je upozorněn s~odkazem na konkrétní záznam, k~jehož údajům je třeba se vrátit.

\subsection{Řízení průjezdové infrastruktury}
\label{sec:zavory}
Stávající systém pro správu závor od společnosti Technopark využívá pro programování ovládací jednotky závor program Kartoun\footnote{Informace o~systému viz \url{https://eshop.technopark.cz/z6987-kartoun-software-pro-spravu-pristupoveho-systemu}.}, jenž běží na lokálním počítači v~areálu.
Tento software není ovládacím prvkem pro samotné otevírání, ale slouží jako správce autonomní paměti vjezdových jednotek.
Kartoun využívá jako své úložiště databázi ve formátu \texttt{.mdb}, jež obsahuje whitelist (seznam povolených) RFID karet.

Hlavní funkcí programu Kartoun je kontinuální synchronizace: jakmile se v~jeho databázi objeví nový identifikátor karty, program jej automaticky nahraje do vnitřní nevolatilní paměti fyzické ovládací jednotky závory.
Závora pak při přiložení karty nečeká na pokyn z~počítače, ale sama porovná ID karty se svou vnitřní pamětí a~otevře se.
Tento mechanismus zajišťuje, že vjezd funguje i~při vypnutém počítači nebo výpadku sítě, pokud je karta již v~jednotce uložena.

Protože software Kartoun je uzavřený proprietární systém a~nenabízí žádné programové rozhraní, je nutné se na něj napojit skrze jeho jediné přístupné místo\,--\,zmíněnou databázi \texttt{.mdb}.
Tato databáze se tak stává integračním bodem, jehož prostřednictvím spolu oba systémy mohou komunikovat.
Vzhledem k~potřebě propojit backend s~tímto lokálním systémem je na počítači přímo připojeném k~ovládací jednotce závor provozován lehký REST API server, jenž slouží jako komunikační most.
Tento server přijímá instrukce z~backendu a~následně provádí zápisy do sdílené databáze \texttt{.mdb}, již Kartoun kontinuálně monitoruje.
Vzhledem k~tomu, že databáze \texttt{.mdb} je sdílena dvěma procesy bez nativního zámkového mechanismu, je přístup k souboru z~obou stran chráněn na úrovni souborového systému; v~případě selhání zápisu integrační vrstva pokus opakuje s~exponenciálním odstupem a~v~krajním případě eskaluje stav obsluze.
Kartoun v~tomto schématu figuruje jako nezbytný komunikační most mezi novým systémem a~fyzickou ovládací jednotkou.

\section{Uživatelské rozhraní}
Proces návrhu uživatelského rozhraní (UI) vychází z~analýzy klíčových potřeb uživatelů, jež je následně transformována do podoby strukturních schémat, takzvaných wireframů.
Tyto nízkoúrovňové modely slouží k~definici informační architektury, rozvržení ovládacích prvků a~celkové navigace, přičemž záměrně abstrahují od vizuálního stylu, barevnosti či detailní grafiky.
Cílem je vytvořit logickou hierarchii a~zajistit, aby prioritní informace byly pro uživatele okamžitě dostupné.
V~následující části jsou proto představeny wireframy stěžejních částí systému, jež slouží jako technický podklad pro budoucí implementaci.

Rozhraní na obrázku \ref{fig:ui-1} představuje přehledovou obrazovku určenou pro personál recepce a~slouží jako vstupní obrazovka aplikace po zahájení směny.
Levý navigační panel umožňuje přechod mezi moduly aplikace a~v~dolní části identifikuje přihlášeného uživatele s~možností ukončit směnu.
Centrální plocha je strukturována do horizontálních pásem podle informační priority\,--\,přehledových karet s~klíčovými ukazateli dne (obsazenost ubytování, počty hostů a~vozidel v~areálu, zbývající příjezdy a~odjezdy, nezaúčtované tržby), operativní agendy (jména přijíždějících a~odjíždějících hostů, stav úklidu, aktuální stav pokladny), upozornění na faktury po splatnosti a~denního stravovacího plánu, doplněných výhledem na akce a~obsazenost v~následujících sedmi dnech.
Toto uspořádání od souhrnných ukazatelů přes denní agendu k~budoucímu výhledu odpovídá obvyklému postupu obsluhy při převzetí směny.

\begin{figure}[hbt!]
    \centering
    \includegraphics[width=0.8\textwidth]{ui/ui-dashboard.pdf}
    \caption{Návrh UI základní navigace v~aplikaci a~přehledové obrazovky.}
    \label{fig:ui-1}
\end{figure}

Modul ubytovacího plánu, zobrazený na obrázku \ref{fig:ui-2}, je koncipován jako interaktivní časová osa určená pro detailní správu a~přehled ubytovacích kapacit v~čase.
V~horní části rozhraní se nacházejí ovládací prvky pro volbu zobrazení, jež uživateli umožňují přepínat mezi aktivní časovou osou (takzvanou plachtou) a~tabulkovým náhledem, případně rychle přistupovat ke zrušeným rezervacím, skupinovým rezervacím a~rezervacím čekajícím na potvrzení.
Součástí horní lišty je rovněž vyhledávací pole a~tlačítko pro vytvoření nové položky související s~rezervacemi.
Jádro modulu tvoří matice, jejíž vertikální osa kategorizuje ubytovací objekty seskupené podle typu (například bungalovy a~zděné chaty), zatímco horizontální osa reprezentuje kalendářní dny s~vyznačením konkrétních rezervací a~stavu jednotlivých jednotek.
Aktuální den je zvýrazněn svislým pruhem napříč celou maticí, čímž je usnadněna časová orientace.
Vizuální kódování rezervací podle jejich stavu (zaplaceno, skupinová, ze skupiny, potvrzená) je vysvětleno v~legendě pod ovládací lištou a~umožňuje obsluze rozpoznat povahu každé rezervace z~jediného pohledu.
Celková ergonomie návrhu je podpořena spodní navigační lištou pro rychlý časový posun mezi měsíci doplněnou o~přepínač roku, což recepci usnadňuje operativní plánování i~zvyšuje čitelnost obsazenosti areálu v~dlouhodobějším horizontu.

\begin{figure}[hbt!]
    \centering
    \includegraphics[width=0.8\textwidth]{ui/ui-plachta.pdf}
    \caption{Návrh UI s~přehledem rezervací v~rezervační plachtě.}
    \label{fig:ui-2}
\end{figure}

Modul plánování stravování, zobrazený na obrázku \ref{fig:ui-strava}, je koncipován jako provozní nástroj recepce pro plánování objednávek a~časů výdeje stravy v~rámci jednoho týdne.
V~horní části rozhraní se nachází přepínač týdnů doplněný o~rychlý souhrn počtu ubytovaných skupin a~celkového počtu porcí v~aktuálně zvoleném týdnu, čímž obsluha okamžitě vidí míru vytížení.
Hlavní obsah je rozdělen do dvou paralelních panelů sdílejících společnou časovou osu týdne.
První panel zobrazuje pro každou ubytovanou skupinu počty porcí podle jednotlivých jídel\,--\,snídaně, oběda a~večeře\,--\,s~dalším rozkladem na dietní varianty (normální, vegetariánskou, bezlepkovou a~bezlaktózovou), doplněné o~datum příjezdu, odjezdu a~celkový počet osob.
Druhý panel pak pro stejné skupiny organizuje konkrétní časy výdejů jídel a~kontaktní osoby zodpovědné za skupinu; výchozí časy jsou patrné z~hlavičky sloupců a~mohou být pro jednotlivé skupiny individuálně upraveny podle jejich denního programu.
Rozdělení do dvou samostatných panelů odráží skutečnost, že časy výdejů přímo nesouvisí s~objednanými počty porcí; jejich oddělení do samostatných mřížek tak zvyšuje přehlednost obou agend a~umožňuje s~nimi pracovat nezávisle.
Aktuální den je v~obou panelech zvýrazněn svislým pruhem, čímž je recepčním usnadněna orientace v~okamžitě platné agendě.

\begin{figure}[hbt!]
    \centering
    \includegraphics[width=0.8\textwidth]{ui/ui-strava.pdf}
    \caption{Návrh UI plánování objednávek a~výdejů stravy.}
    \label{fig:ui-strava}
\end{figure}

Pohled „Místa“ v~mobilní aplikaci pro provozní personál, zobrazený na obrázku \ref{fig:ui-mista}, slouží pro zobrazení aktuálního stavu všech ubytovacích jednotek v~areálu.
Účelem je poskytnout pracovníkům pohybujícím se v~terénu okamžitou informaci o~tom, kde je očekávaný příjezd, odjezd, kde je třeba úklidu a~které jednotky jsou mimo provoz, bez nutnosti kontaktovat recepci.
V~horní části obrazovky se nachází souhrnná lišta pěti karet s~počty jednotek v~každém ze stavů, jež zároveň plní funkci legendy pro vizuální kódování dlaždic v~mřížce níže.
Jádro pohledu tvoří mřížka ubytovacích jednotek seskupených podle typu, kde každá dlaždice vyjadřuje současně několik údajů: stav místa vizuálním stylem pozadí, označení chaty, datum plánovaného odjezdu a~indikátor potřeby úklidu.
Personál tak dokáže rozpoznat povahu každého místa z~jediného pohledu.
Dolní navigační lišta umožňuje rychlý přechod mezi pěti hlavními moduly mobilní aplikace (úklid, údržba, strava, místa, skener QR kódů) a~její umístění respektuje ergonomii jednoruční obsluhy mobilního zařízení.

\begin{figure}[hbt!]
    \centering
    \includegraphics[width=0.45\textwidth]{ui/ui-mista.pdf}
    \caption{Návrh UI navigace v~mobilní aplikaci a~zobrazení stavu jednotlivých míst.}
    \label{fig:ui-mista}
\end{figure}